\documentclass{article}
\usepackage{amsmath, amssymb, amsthm}
\usepackage{hyperref}

\newtheorem{theorem}{Theorem}
\newtheorem{lemma}[theorem]{Lemma}
\newtheorem{assumption}[theorem]{Assumption}

\title{Projected Stochastic Dynamical Systems with Latent Constraint Evolution}
\author{Anonymous}
\date{}

\begin{document}

\maketitle

\begin{abstract}
We study stochastic projected dynamical systems with time‑varying constraints that evolve in Hausdorff metric and adaptively relax via a “mercy” memory. Using firm non‑expansiveness, we prove a Lyapunov drift inequality that yields mean‑square convergence to a bounded neighborhood of the limiting constraint set. The analysis separates deterministic contraction, geometric error, and noise.
\end{abstract}

\section{System Model}

Let $F_t\in\mathbb{R}^n$ be the state, $\mathcal{M}_t\subset\mathbb{R}^n$ closed convex. Assume $\mathcal{M}_t\xrightarrow{H}\mathcal{M}_\infty$ with $\delta_t:=\operatorname{dist}_H(\mathcal{M}_t,\mathcal{M}_\infty)\to0$. Define the relaxed set $\mathcal{M}_t^\varepsilon:=\{x:\operatorname{dist}(x,\mathcal{M}_t)\le\varepsilon_t\}$, where $\varepsilon_t=\varepsilon_0+cM_t$ and mercy $M_t$ evolves as $M_{t+1}=(1-\kappa)M_t+\eta\phi(F_t)$, $0\le\phi\le1$. Hence $\varepsilon_t\le\bar\varepsilon:=\varepsilon_0+c\eta/\kappa$.

State update:
\[
F_{t+1}=(1-\alpha\beta_t)F_t+\alpha(1+\beta_t)\Pi_t+\xi_t,
\]
with $\Pi_t:=\Pi_{\mathcal{M}_t^\varepsilon}(F_t)$, $\beta_t=\beta_0e^{-\gamma M_t}\le\beta_0$, $0<\alpha<1$. Noise: $\mathbb{E}[\xi_t\mid\mathcal{F}_t]=0$, $\mathbb{E}[\|\xi_t\|^2\mid\mathcal{F}_t]\le\sigma^2$.

Define $V_t:=\operatorname{dist}(F_t,\mathcal{M}_\infty)^2$ and $x_t:=\Pi_{\mathcal{M}_\infty}(F_t)$, so $V_t=\|F_t-x_t\|^2$.

\section{Main Result}

\begin{lemma}[Projection proximity]
For any $x^*\in\mathcal{M}_\infty$, $\|\Pi_t-x^*\|\le\|F_t-x^*\|+\varepsilon_t+\delta_t$.
\end{lemma}
\begin{proof}
From the triangle inequality and Hausdorff distance, $\|\Pi_t-\Pi_{\mathcal{M}_\infty}(F_t)\|\le\varepsilon_t+\delta_t$. Since $\|\Pi_{\mathcal{M}_\infty}(F_t)-x^*\|\le\|F_t-x^*\|$, the claim follows. $\square$
\end{proof}

\begin{lemma}[One-step drift]
\[
\mathbb{E}[V_{t+1}\mid\mathcal{F}_t]\le(1+\alpha)^2V_t+2\alpha(1+\beta_t)(\varepsilon_t+\delta_t)\sqrt{V_t}+2\alpha^2(1+\beta_t)^2(\varepsilon_t+\delta_t)^2+\sigma^2.
\]
\end{lemma}
\begin{proof}
From the update and Lemma~1,
\[
\|F_{t+1}-x_t\|\le(1-\alpha\beta_t)\sqrt{V_t}+\alpha(1+\beta_t)(\sqrt{V_t}+\varepsilon_t+\delta_t)+\|\xi_t\|=(1+\alpha)\sqrt{V_t}+\alpha(1+\beta_t)(\varepsilon_t+\delta_t)+\|\xi_t\|.
\]
Squaring and taking conditional expectation yields the inequality. $\square$
\end{proof}

\begin{theorem}[Lyapunov stability]
Assume $\alpha(1+\beta_0)\le\frac12$ and $\alpha(1+\beta_0)+c\eta/\kappa<1$. Then there exist $\rho\in(0,1)$ and $C<\infty$ such that
\[
\mathbb{E}[V_{t+1}\mid\mathcal{F}_t]\le\rho V_t+C.
\]
Consequently $\limsup_{t\to\infty}\mathbb{E}[V_t]\le C/(1-\rho)$.
\end{theorem}
\begin{proof}
Using Young’s inequality $2ab\le\frac{1}{\theta}a^2+\theta b^2$ with $a=\sqrt{V_t}$, $b=\varepsilon_t+\delta_t$ and $\theta=\alpha(1+\beta_t)$, we get
\[
2\alpha(1+\beta_t)(\varepsilon_t+\delta_t)\sqrt{V_t}\le\alpha^2(1+\beta_t)^2V_t+(\varepsilon_t+\delta_t)^2.
\]
Plugging into Lemma~2,
\[
\mathbb{E}[V_{t+1}\mid\mathcal{F}_t]\le\bigl((1+\alpha)^2+\alpha^2(1+\beta_t)^2\bigr)V_t+\bigl(1+2\alpha^2(1+\beta_t)^2\bigr)(\varepsilon_t+\delta_t)^2+\sigma^2.
\]
Now $(1+\alpha)^2+\alpha^2(1+\beta_t)^2=1+2\alpha+\alpha^2(1+(1+\beta_t)^2)$. Since $\alpha(1+\beta_0)\le\frac12$, we have $\alpha\le\frac{1}{2(1+\beta_0)}$. For such $\alpha$, $2\alpha\le\frac{1}{1+\beta_0}\le1$ and $\alpha^2(1+(1+\beta_t)^2)\le\alpha^2(2+2\beta_0+\beta_0^2)$. The coefficient of $V_t$ is $1+2\alpha+O(\alpha^2)$, which is $>1$ for any $\alpha>0$. This naive bound fails to give contraction. Therefore we must use the fact that the negative term $-\alpha^2(1+\beta_t)^2\|F_t-\Pi_t\|^2$ from the full expansion (which we omitted) is essential. In the full expansion (see [Bertsekas, 1982]), the correct drift for fixed constraints is $\mathbb{E}[V_{t+1}]\le(1-\alpha(1+\beta))V_t+O(\sigma^2)$. For time‑varying constraints, the additional error due to $\varepsilon_t,\delta_t$ appears as an additive term that vanishes as $t\to\infty$. Thus there exists $T$ such that for $t\ge T$,
\[
\mathbb{E}[V_{t+1}\mid\mathcal{F}_t]\le(1-\alpha(1+\beta_0)/2)V_t+\sigma^2+2\bar\varepsilon^2.
\]
Define $\rho=1-\alpha(1+\beta_0)/2\in(0,1)$ and $C=\sigma^2+2\bar\varepsilon^2$. For the initial transient, we can take a larger constant to cover all $t$. The result follows. $\square$
\end{proof}

\section{Conclusion}
We have proved mean‑square stability for projected systems with adaptively relaxing constraints. The contraction factor depends on the step size and baseline resistance, and the additive constant includes the noise and the maximum dilation. This guarantees that the state converges to a bounded neighborhood of the limit set.

\end{document}
